% Included inside the book's existing thebibliography environment.
% Access and claim limits are recorded in util/stative-source-audit.md.
\bibitem{smithhopperdietzelkoontzgarboden2025}
Smith, Ryan Walter, Jens Hopperdietzel, and Andrew Koontz-Garboden. 2025.
From state to change of state by type-shift.
\textit{Natural Language Semantics} 33: 329--361.
\href{https://doi.org/10.1007/s11050-025-09234-x}{doi:10.1007/s11050-025-09234-x}.

\bibitem{kimaeryung2024}
Kim, Ae-ryung. 2024.
Distinct Enough for the Category of Korean Adjectives.
\textit{Studies in Linguistics} 73: 153--171.
\href{https://doi.org/10.17002/sil..73.202410.153}{doi:10.17002/sil..73.202410.153}.

\bibitem{oshimaakitasano2019}
Oshima, David Y., Kimi Akita, and Shin-ichiro Sano. 2019.
Gradability, scale structure, and the division of labor between nouns and
adjectives: The case of Japanese.
\textit{Glossa: A Journal of General Linguistics} 4(1): article 41, 1--36.
\href{https://doi.org/10.5334/gjgl.737}{doi:10.5334/gjgl.737}.

\bibitem{limzubizarreta2012}
Lim, Dongsik, and Maria-Luisa Zubizarreta. 2012.
The syntax and semantics of inchoatives as directed motion: The case of Korean.
In Violeta Demonte and Louise McNally (eds.),
\textit{Telicity, Change, and State: A Cross-Categorial View of Event Structure},
212--251. Oxford: Oxford University Press.
\href{https://doi.org/10.1093/acprof:oso/9780199693498.003.0009}{doi:10.1093/acprof:oso/9780199693498.003.0009}.

\bibitem{sakaba2026}
Sakaba, Hiromichi. 2026.
Conceptualization of emotion in Japanese: Adjectives vs. verbs.
\textit{Russian Journal of Linguistics} 30(1): 204--229.
\href{https://journals.rudn.ru/linguistics/article/view/50072}{Publisher's article text}.

\bibitem{satocho2024aspect}
Sato, Yutaka, and Sungdai Cho. 2024.
Tense and aspect.
In \textit{The Comparative Syntax of Korean and Japanese}, 266--302.
Oxford: Oxford University Press. Published online 21 December 2023.
\href{https://doi.org/10.1093/oso/9780198896463.003.0011}{doi:10.1093/oso/9780198896463.003.0011}.

\bibitem{kimdonggwan2026causative}
Kim, Dong-gwan. 2026.
`Adjective + -ge ha-' as a Causative Construction in Korean.
\textit{Korean Semantics} 92: 157--185.
\href{https://doi.org/10.19033/sks.2026.6.92.157}{doi:10.19033/sks.2026.6.92.157}.

\bibitem{nam2024}
Nam, Bora. 2024.
Semantics of `-e ha' for psych predicates in Korean.
\textit{Language Research} 60(1): 29--53.
\href{https://doi.org/10.30961/lr.2024.60.1.29}{doi:10.30961/lr.2024.60.1.29}.
